\chapter{Conclusion and Future Work}\label{ch:conclusion}

\section{Summary of Contributions}\label{sec:summary_contributions}

This thesis began with the question of how to inspect and question the evidence used by a trained vision model. When attribution informs model inspection, correction, or intervention, understanding the calculation behind an explanation becomes essential. The thesis develops a mathematical framework for this purpose: it exposes opposing contributions within trained networks, characterizes backward attribution through games, and evaluates methods derived from these representations.

The two papers establish complementary benefits. \emph{Hidden Monotonicity} makes structured split-streams available for attribution without retraining or changing the network's predictions~\parencite{hiddenMono}. \emph{Playing the Network Backward} provides a common mathematical language for comparing and adapting gradient and LRP rules~\parencite{playingBackward2026}. The thesis connects these perspectives and extends the decomposition framework to smooth activations and Transformers.

\paragraph{What structure can be exposed without changing the network's function?}
We have presented a common framework for DC, stopping, mixing, and regularized decompositions. For pretrained ReLU networks and CNNs, the DC construction exposes monotone and convex split-streams on doubled inputs whose difference recovers the original function. Stabilized forward and backward passes make this representation usable for attribution. The absolute-curvature construction supplies explicit splits for GELU, Swish, sigmoid, and Softplus, with scalar profiles of minimal curvature. Convex product, softmax, and LayerNorm pairs yield an exact DC decomposition of arbitrary finite Transformers of the stated architecture (Section~\ref{sec:transformer_dc}). These results establish that useful mathematical structure can be exposed in an unconstrained model while preserving its predictions.

\paragraph{What does a backward attribution actually calculate?}
The Stopping and Routing Games represent gradient and $\alpha\beta$-\ac{lrp} attribution as equilibrium quantities on per-neuron states. At $\tau=1$, the active player's equilibrium values at Routing Game activation states recover the network activations in log space. Under the stated assumptions, signed discounted occupation measures recover gradients in the Stopping Game and $\alpha\beta$-\ac{lrp} in the Routing Game.

Beyond these recovery results, the framework introduces a new language for describing and designing attribution methods. Temperature expresses routing concentration, risk aversion specifies how uncertain continuation values are treated, and policy regularization, extended action sets, and smooth gates describe further choices in the backward calculation. Desired explanation properties can thus be formulated as explicit changes to objectives, policies, and available actions, providing a common basis for comparing methods and deriving new rules.

Equilibrium policies induce probability distributions over backward trajectories, from which attribution maps are obtained through signed discounted projection. For trajectory laws on aligned game graphs, their Hellinger distance is computable in one backward pass and compares routing before this projection. The randomization experiments demonstrate its practical distinction from map comparisons: trajectory distributions change even when projected heatmaps remain similar. Attribution can therefore be examined both through the resulting map and through the modeled routing that produces it.

\paragraph{How can these representations guide the design and evaluation of attribution methods for visual inspection?}
SplitCAM, SplitLRP, SplitGrad, and the game variants demonstrate the practical value of these representations. Consistent with the conclusions of \emph{Hidden Monotonicity}, the CNN experiments support splitting as a tool for improving several attribution criteria on pretrained networks. SplitCAM has the best observed localization and Selectivity means on VGG-16 and the lowest observed Selectivity mean on ResNet-18, where it ties with the best baseline on the Pointing Game, while SplitLRP has the lowest observed Maximum Sensitivity in the VGG-16 comparison. These comparisons are descriptive: the split campaign reports means without uncertainty estimates, and small differences, such as the ResNet-18 Selectivity gap of $0.014$ to Guided Backprop, have not been assessed. SplitLRP's nearly invariant ranks under classifier randomization (Table~\ref{tab:randomization_sanity_check}) limit interpreting its stability as model-specific explanation quality. Different configurations lead on different criteria.

The game experiments demonstrate that principled changes of backward routing can produce strong attribution methods. A selected Routing Game configuration achieves Attribution Localization of $0.711$ on ViT-B/16, compared with $0.644$ for the best baseline in the original comparison. The localization advantage transfers to Pascal VOC without retuning the game parameters. As emphasized in \emph{Playing the Network Backward}, these gains coexist with stronger faithfulness results for LRP rules. The controls clarify the distinction between concentrating attribution mass, ranking relevant features, and obtaining robust explanations: the routing modes give controllable, strong localization under the stated protocol, and this localization lead does not come with a corresponding advantage in perturbation faithfulness.

The additional ViT DC experiments use the explicit GELU split and raise the reported CAM localization means at selected layer readouts with fixed attention and LayerNorm (Section~\ref{sec:vit-extension}). The Pointing Game gain over LayerCAM is significant after Holm correction, whereas the adjusted tests do not resolve the Attribution Localization and Top-$k$ differences. Directly trained difference models provide a complementary architectural perspective: their MNIST split-stream gradients exhibit class-dependent sensitivities to present strokes and candidate missing-stroke locations. Together, these experiments show how structural representations can support attribution design and generate hypotheses about model behavior.

\section{Limitations and Broader Impact}\label{sec:limitations}\label{sec:broader_impact}\label{ptnb:app:broader-impact}

The evidence concerns the tested models, datasets, and protocols. Human-mask agreement and decision faithfulness are assessed by complementary metrics. The game identities for ViT condition on attention, LayerNorm statistics, and GELU gates. Hellinger distance compares the modeled trajectory-probability laws, rather than every accompanying discount, value, or operation detached from the backward calculation. Its current applications demonstrate sensitivity to routing changes, without establishing that the distance predicts explanation quality.

The exact structural guarantees concern the split-streams; input-dependent stabilization need not preserve their global convexity or monotonicity. Exact Transformer identities also do not prevent floating-point cancellation. Stabilization and layer selection remain consequential practical choices. Attribution can support scrutiny of consequential systems, but favorable metrics and visually convincing maps alone cannot justify trust in their behavior.

\section{Future Research Directions}\label{sec:future_directions}

\paragraph{Concepts, interventions, and model correction.}
Concept-conditioned trajectory events, informed by \ac{nrm} queries (Section~\ref{sec:related_work_games}), could connect backward routing to interpretable features and internal components. A central next step is to test whether these diagnostics predict the effects of interventions on the complete model. Dedicated benchmarks could compare intervention targets identified through split-streams, trajectory information, and existing attribution maps, measuring both changes in unwanted dependencies and preservation of desired behavior. This would connect the mathematical representations to the model-inspection and correction tasks motivating the thesis. Trajectory distances could also enter training directly: a Hellinger penalty between the routing laws of background-swapped pairs, as in Waterbirds~\parencite{Sagawa2020GroupDRO}, or between the previous and current model on replayed inputs in continual learning, would constrain how a prediction is computed rather than only its map or output~\parencite{Fel2026MechanisticShaping}. Such penalties cover only the components modeled in the Markov process and lose gradient where the distance saturates near one.

\paragraph{Comparing targets.}
Chapter~\ref{ch:trajectory} compares trajectory laws with a shared target under input noise and parameter randomization. For a fixed image, \textcite{Maess2026Spectral} stack normalized attribution maps for different target outputs as matrix columns and use its singular values to measure whether targets yield distinct maps. The same comparison could be applied to trajectory laws of different targets, before projection onto input pixels.

\paragraph{From individual calculations to broader behavioral patterns.}
Hellinger distance may also become useful independently of heatmap evaluation. Clustering inputs by distances between their modeled trajectory laws could reveal groups of examples processed through similar routes. Examining representative inputs and their local explanations could provide insight linking local explanations to global patterns of network behavior.

A more exploratory application would cache reference trajectory laws of a language model under controlled truthful and deliberately misleading conditions, then compare subsequent computations with these references as a candidate \emph{lie detector}. Such experiments could supply the correct answer while inducing the model to misreport it. They would require comparable trajectory representations and tests separating deception from topic, wording, uncertainty, and ordinary mistakes. Both applications remain untested: they motivate further study of distances between modeled calculations without assuming that routing similarity already identifies a semantic or behavioral property.

\paragraph{Choosing decompositions, games, and layers.}
Neither a network's function nor its architecture uniquely determines an explanatory representation. Adding a common admissible convex component to both split-streams preserves their difference while potentially changing split-stream-specific attribution; different games, policies, and discounts can likewise expose different aspects of a computation. Future work should identify which explanatory properties are invariant under such choices and develop selection principles tied to the intended task. The strong dependence on layer selection motivates automatic selection and aggregation methods that generalize beyond validation-based tuning on one dataset.

\paragraph{Directly trained opposing split-stream architectures.}
The MNIST results motivate a broader investigation of training models as differences of structured networks. Larger datasets and architectures should test whether useful split-stream roles persist across training runs and whether architectural or regularization choices can encourage them without unacceptable losses in predictive performance. Complete-model interventions are needed to assess the proposed present-feature and missing-feature interpretations and their usefulness for model correction.

\paragraph{Robustness through trajectory distances.}
Relating perturbations of trajectory laws to changes in signed discounted occupations could yield robustness guarantees for attribution. Such bounds must control the magnitude and variation of the path weights and discounts, as well as changes in the projection when models or rules differ. This would clarify which aspects of explanation stability Hellinger distance can characterize and how these guarantees compare with curvature- or Lipschitz-based bounds.

\paragraph{Numerical stability and practical Transformer decomposition.}
The Transformer construction primarily establishes the mathematical possibility of the split. Limited full-model reconstruction checks and attribution experiments at selected later-layer readouts do not yet establish effective DC attribution throughout the full Transformer. The practical attribution study uses fixed attention and LayerNorm, while exploiting the complete decomposition requires retaining their input dependence. New stabilization strategies are required to control split-stream growth and cancellation while preserving useful individual split-stream signals across attention products, normalization, and depth. Forward or gradient reconstruction alone is insufficient to establish this.

These directions return to the thesis's motivating aim: to understand and question the evidence used by trained networks. The thesis supplies both structural representations and a language for describing and designing backward attribution, making parts of the computation explicit and experimentally accessible. Their further value will depend on whether this access helps us identify meaningful patterns, predict intervention outcomes, and improve model behavior.